\documentclass[conference]{IEEEtran}
\IEEEoverridecommandlockouts

\usepackage{cite}
\usepackage{amsmath,amssymb,amsfonts}
\usepackage{algorithmic}
\usepackage{graphicx}
\usepackage{textcomp}
\usepackage{xcolor}
\usepackage{multirow}
\usepackage{booktabs}
\usepackage{hyperref}
\usepackage{algorithm}

\title{
An Expert-Guided Multi-Modal AI Ecosystem for\\
Diagnostic Intelligence
}

\author{
\IEEEauthorblockN{
Farhan Kashif\IEEEauthorrefmark{1},
Ahmed Sultan\IEEEauthorrefmark{1},
Imama\IEEEauthorrefmark{1}
}
\IEEEauthorblockA{
\IEEEauthorrefmark{1}
School of Electrical Engineering and Computer Science\\
National University of Sciences and Technology (NUST), Islamabad, Pakistan\\
Email: \{fkashif, asultan, imama\}.bese22seecs@seecs.edu.pk
}
}

\begin{document}
\maketitle

\begin{abstract}
Accurate brain tumor segmentation and interpretable clinical reporting are critical for effective diagnosis and treatment planning. While deep learning has achieved strong performance on medical imaging tasks, existing systems remain static, lack adaptability to evolving datasets, and fail to produce clinically actionable outputs. This paper presents the architectural design of an expert-guided multimodal framework that integrates continual learning-based segmentation, atlas-driven anatomical mapping, and evidence-grounded report generation. We extend the Mixture-of-Modality-Experts (MoME+) architecture with Elastic Weight Consolidation (EWC) and experience replay mechanisms designed to enable sequential learning across heterogeneous MRI datasets without catastrophic forgetting. Segmentation outputs are designed to be mapped to standardized brain atlases and encoded into structured JSON schemas containing lesion types, volumetric measurements, and anatomical locations. These schemas serve as factual inputs to a proposed MedGemma-4B medical language model using parameter-efficient LoRA adaptation for generating radiology-style reports with explicit provenance. We implement and evaluate a baseline MoME+ segmentation model on the BraTS 2021 dataset and conduct preliminary continual learning experiments, demonstrating the viability of the proposed architecture. Complete empirical validation across multiple datasets and full LLM integration remain as future work.
\end{abstract}

\begin{IEEEkeywords}
Brain tumor segmentation, multimodal MRI, continual learning, mixture of experts, medical language models, structured report generation, explainable AI
\end{IEEEkeywords}

\section{Introduction}

Gliomas account for approximately 80\% of all malignant brain tumors and present significant diagnostic challenges due to their heterogeneous appearance across imaging modalities \cite{brats2021}. Accurate segmentation of tumor sub-regions—including enhancing tumor (ET), necrotic core (NCR), and peritumoral edema—is essential for treatment planning, yet manual delineation remains time-consuming and subject to inter-observer variability \cite{menze2015brats}.

Deep learning has revolutionized medical image segmentation, with U-Net variants and transformer-based architectures achieving state-of-the-art performance on benchmark datasets \cite{isensee2021nnu, hatamizadeh2022swin}. However, these models face critical limitations in clinical deployment: (1) they are trained on static datasets and cannot adapt to new imaging protocols or pathologies without retraining, (2) they lack robustness to missing or corrupted modalities, and (3) they produce pixel-level outputs without clinical contextualization.

Recent work on Mixture-of-Modality-Experts (MoME) \cite{zhang2024mome} addresses modality-specific feature extraction through specialized expert networks, demonstrating improved robustness to heterogeneous inputs. Concurrently, medical large language models (LLMs) like Med-PaLM \cite{singhal2023medpalm} have shown promise in clinical text generation, though direct vision-language approaches risk hallucination and lack transparency.

This work addresses these limitations through three key contributions:

\begin{enumerate}
\item \textbf{Continual Learning-Enabled Segmentation:} We extend the MoME+ architecture with Elastic Weight Consolidation (EWC) and experience replay to enable sequential learning across multiple datasets without catastrophic forgetting, adapting to evolving clinical data distributions.

\item \textbf{Atlas-Based Structured Representation:} We develop a deterministic pipeline that maps segmentation outputs to standardized brain atlases (Harvard-Oxford), encoding anatomical information into a JSON schema that serves as an explicit, auditable knowledge source for downstream reasoning.

\item \textbf{Evidence-Grounded Report Generation:} We fine-tune MedGemma-4B using parameter-efficient LoRA adaptation on synthetic JSON-to-report pairs, ensuring that generated radiology reports are factually grounded in segmentation outputs with sentence-level provenance.
\end{enumerate}

The proposed framework integrates these components into a unified system that learns continuously, produces interpretable outputs, and maintains clinical trust through explicit fact-checking mechanisms.

\section{Related Work}

\subsection{Multimodal Brain Tumor Segmentation}

\subsubsection{CNN-Based Architectures}

U-Net \cite{ronneberger2015unet} revolutionized medical image segmentation through its encoder-decoder architecture with skip connections, enabling precise localization while maintaining contextual information. The architecture's symmetric design allows gradual downsampling in the encoder (capturing semantic features) and upsampling in the decoder (recovering spatial resolution), with skip connections fusing multi-scale features.

The 3D extension of U-Net processes volumetric medical images directly, avoiding information loss from 2D slice-based approaches. However, 3D convolutions dramatically increase computational cost: a 3×3×3 kernel has 27 parameters versus 9 for 2D, and memory requirements scale cubically with volume dimensions. This necessitates trade-offs between architectural depth and input resolution.

The nnU-Net framework \cite{isensee2021nnu} automated these design choices through a self-configuring pipeline that adapts preprocessing, network architecture, and training strategies to dataset characteristics. Key innovations include: (1) automatic determination of patch size based on GPU memory and median image dimensions, (2) adaptive batch size calculation ensuring stable gradient estimates, (3) dataset-specific data augmentation intensity, and (4) ensemble prediction combining multiple model configurations. nnU-Net achieved top performance on BraTS by optimizing the full segmentation pipeline rather than solely focusing on architectural novelty.

\subsubsection{Transformer-Based Architectures}

More recently, transformer-based architectures have leveraged self-attention mechanisms for improved global context modeling. UNETR \cite{hatamizadeh2022unetr} introduced Vision Transformers (ViT) to 3D medical segmentation by treating volumetric patches as tokens. The encoder applies multi-head self-attention:

\begin{equation}
\text{Attention}(Q, K, V) = \text{softmax}\left(\frac{QK^T}{\sqrt{d_k}}\right)V
\end{equation}

where $Q$, $K$, $V$ are query, key, and value matrices derived from input patches, and $d_k$ is the key dimension. This enables every patch to attend to all other patches, capturing long-range dependencies that CNNs struggle with due to limited receptive fields.

Swin UNETR \cite{hatamizadeh2022swin} improved computational efficiency through shifted window attention, restricting self-attention to local windows while enabling inter-window communication via window shifting across layers. This reduces complexity from $O(n^2)$ to $O(n)$ where $n$ is the number of patches, making transformer models practical for high-resolution 3D volumes.

\subsubsection{Attention Mechanisms in CNNs}

Several attention mechanisms enhance CNN feature representations without the computational burden of full transformers. Squeeze-and-Excitation (SE) blocks \cite{hu2018squeeze} apply channel-wise attention:

\begin{equation}
\mathbf{z} = F_{sq}(\mathbf{U}) = \frac{1}{H \times W \times D} \sum_{i,j,k} \mathbf{U}_{i,j,k}
\end{equation}

\begin{equation}
\mathbf{s} = F_{ex}(\mathbf{z}, W) = \sigma(W_2 \delta(W_1 \mathbf{z}))
\end{equation}

where $F_{sq}$ performs global average pooling, $F_{ex}$ applies two fully connected layers with ReLU ($\delta$) and sigmoid ($\sigma$) activations, and $\mathbf{s}$ represents channel-wise scaling weights. This adaptively recalibrates feature responses based on channel importance.

Convolutional Block Attention Module (CBAM) \cite{woo2018cbam} extends this with spatial attention, first computing channel attention followed by spatial attention that identifies important regions within feature maps. The sequential application allows the model to focus on both "what" (channels) and "where" (spatial locations) to emphasize.

\subsection{Mixture of Experts for Multimodal Learning}

\subsubsection{Modality-Specific Processing}

Multi-modal medical imaging presents unique challenges: different modalities (T1, T1ce, T2, FLAIR) capture distinct tissue properties, have varying contrast mechanisms, and exhibit different noise characteristics. Traditional early fusion (concatenating modalities as input channels) or late fusion (combining prediction from separate models) strategies fail to exploit modality-specific patterns optimally.

Zhang et al. \cite{zhang2024mome} proposed the Mixture-of-Modality-Experts (MoME) framework, assigning dedicated expert networks to each modality. Each expert $E_m$ processes only its corresponding modality:

\begin{equation}
\mathbf{f}_m = E_m(\mathbf{x}_m; \theta_m)
\end{equation}

where $\mathbf{x}_m$ is the input for modality $m$ and $\theta_m$ are modality-specific parameters. This specialization allows each expert to learn representations tailored to its input modality's characteristics without interference from other modalities during feature extraction.

\subsubsection{Gating Networks for Expert Fusion}

The gating network $G$ dynamically weights expert contributions based on input characteristics:

\begin{equation}
\mathbf{w} = \text{softmax}(G(\mathbf{x}; \phi))
\end{equation}

where $\phi$ are gating network parameters and $\mathbf{w} \in \mathbb{R}^M$ contains expert weights summing to 1. The softmax ensures a valid probability distribution over experts. Key design choices include:

\begin{itemize}
\item \textbf{Input representation:} The gating network can receive raw inputs, downsampled inputs, or intermediate features
\item \textbf{Architecture depth:} Shallow networks (1-2 layers) for simple weighting vs. deep networks for complex routing
\item \textbf{Temperature scaling:} A learnable temperature parameter $\tau$ controls the sharpness of the softmax distribution
\end{itemize}

MoME demonstrated 2-4\% Dice improvement over nnU-Net across nine brain lesion datasets, showing that modality-specific feature extraction with learned fusion outperforms naive concatenation strategies.

\subsection{Continual Learning in Medical Imaging}

\subsubsection{Catastrophic Forgetting Problem}

Continual learning (CL) enables models to learn sequentially from non-stationary data distributions while preserving previously acquired knowledge \cite{parisi2019continual}. The central challenge is \textit{catastrophic forgetting}: when neural networks are trained on new tasks, gradient descent updates can dramatically degrade performance on previous tasks by overwriting previously learned parameters.

Formally, given a sequence of tasks $\{T_1, T_2, \ldots, T_K\}$ with datasets $\{\mathcal{D}_1, \mathcal{D}_2, \ldots, \mathcal{D}_K\}$, the objective is to minimize:

\begin{equation}
\min_{\theta} \sum_{k=1}^{K} \mathbb{E}_{(\mathbf{x}, \mathbf{y}) \sim \mathcal{D}_k}[\mathcal{L}_k(f(\mathbf{x}; \theta), \mathbf{y})]
\end{equation}

while maintaining competitive performance on all tasks after training on task $K$, not just the current task.

\subsubsection{Elastic Weight Consolidation}

Elastic Weight Consolidation (EWC) \cite{kirkpatrick2017ewc} addresses catastrophic forgetting through regularization based on parameter importance. The core insight is that not all parameters contribute equally to previous task performance—some are critical while others are less important.

EWC estimates parameter importance using the Fisher Information Matrix (FIM), which measures the sensitivity of the loss function to parameter perturbations:

\begin{equation}
F_i = \mathbb{E}_{(\mathbf{x},\mathbf{y}) \sim \mathcal{D}_k}\left[\left(\frac{\partial \log p(\mathbf{y}|\mathbf{x}; \theta)}{\partial \theta_i}\right)^2\right]
\end{equation}

High Fisher values indicate parameters that strongly affect predictions on task $k$—changing them significantly would degrade task $k$ performance. The EWC loss penalizes deviations from optimal previous task parameters weighted by their importance:

\begin{equation}
\mathcal{L}_{\text{EWC}} = \mathcal{L}_{\text{current}} + \sum_{k=1}^{K-1} \sum_i \frac{\lambda}{2} F_i^{(k)} (\theta_i - \theta_i^{*(k)})^2
\end{equation}

where $\theta_i^{*(k)}$ are optimal parameters after training on task $k$, and $\lambda$ controls regularization strength. This quadratic penalty allows important parameters to change minimally while permitting greater plasticity in less critical parameters.

\subsubsection{Experience Replay}

Experience replay methods \cite{rebuffi2017icarl} maintain a memory buffer $\mathcal{B}$ storing representative samples from previous tasks. During training on task $k$, mini-batches combine current task data with replayed samples:

\begin{equation}
\mathcal{L}_{\text{total}} = \mathcal{L}(\mathcal{D}_k) + \alpha \mathcal{L}(\text{Sample}(\mathcal{B}))
\end{equation}

where $\alpha$ balances current and historical learning. Key design considerations include:

\begin{itemize}
\item \textbf{Buffer management:} Fixed-size buffers require sample selection strategies (random, herding, reservoir sampling)
\item \textbf{Sampling strategy:} Uniform sampling vs. importance sampling based on loss or uncertainty
\item \textbf{Replay frequency:} How many replay samples per current task batch
\end{itemize}

Perkonigg et al. \cite{perkonigg2021continual} demonstrated that hybrid approaches combining EWC and replay achieve superior performance on medical imaging tasks under domain shift, as regularization alone cannot prevent forgetting when task distributions are highly divergent.

\subsection{Medical Report Generation}

\subsubsection{Vision-Language Models}

Automated medical report generation aims to convert imaging findings into natural language descriptions. Early approaches used template-based systems with rule-based text generation, providing high factual accuracy but limited expressiveness and flexibility.

Vision-language models such as R2Gen \cite{chen2020r2gen} and RadGraph \cite{jain2021radgraph} employ encoder-decoder architectures where a CNN encoder processes images and an RNN/Transformer decoder generates reports token-by-token. However, these end-to-end approaches suffer from \textit{hallucination}—generating plausible but factually incorrect statements not supported by the image.

\subsubsection{Grounded Report Generation}

Grounded approaches \cite{zhang2023radiology} introduce intermediate structured representations between vision and language. These typically involve:

\begin{enumerate}
\item \textbf{Fact extraction:} Converting image features into structured facts (e.g., "lesion volume: 27.3 cm³", "location: frontal lobe")
\item \textbf{Fact verification:} Checking that generated text aligns with extracted facts
\item \textbf{Constrained generation:} Conditioning language model on verified facts
\end{enumerate}

This factorization provides explicit provenance—every statement in the report can be traced to a specific extracted fact, enabling human verification and debugging.

\subsubsection{Parameter-Efficient Fine-Tuning}

Medical LLMs including Med-PaLM \cite{singhal2023medpalm}, MedAlpaca \cite{han2023medalp}, and MedGemma \cite{medgemma2024}demonstrate strong performance on medical question-answering. However, full fine-tuning requires updating billions of parameters, demanding substantial computational resources.

Low-Rank Adaptation (LoRA) \cite{hu2022lora} enables efficient adaptation by decomposing weight updates into low-rank matrices. For a pre-trained weight matrix $W_0 \in \mathbb{R}^{d \times k}$:

\begin{equation}
h = W_0 x + \Delta W x = W_0 x + BA x
\end{equation}

where $B \in \mathbb{R}^{d \times r}$ and $A \in \mathbb{R}^{r \times k}$ with rank $r \ll \min(d, k)$. Freezing $W_0$ and training only $A$ and $B$ reduces trainable parameters by orders of magnitude while maintaining performance comparable to full fine-tuning. For a 4B parameter model with $r=16$, LoRA typically trains <0.5\% of total parameters.

\section{Methodology}

\subsection{Overall Architecture}

The proposed framework consists of four tightly integrated stages:

\begin{enumerate}
\item \textbf{Continual Learning-Based Segmentation:} MoME+ architecture extended with EWC and experience replay
\item \textbf{Atlas-Based Anatomical Mapping:} Registration to standard brain atlases for region identification
\item \textbf{Structured JSON Schema Generation:} Deterministic encoding of segmentation and anatomical features
\item \textbf{LLM-Based Report Generation:} Fine-tuned MedGemma-4B with LoRA adapters
\end{enumerate}

\subsection{MoME+ Segmentation with Continual Learning}

\subsubsection{Modality-Specific Expert Networks}

The base MoME+ architecture consists of $N=4$ modality-specific expert networks $\{E_1, E_2, E_3, E_4\}$ corresponding to T1, T1ce, T2, and FLAIR modalities respectively. Each expert $E_i$ follows a 3D U-Net encoder-decoder architecture with residual connections and CBAM attention blocks.

Formally, let $\mathbf{x}_i \in \mathbb{R}^{1 \times D \times H \times W}$ denote the input volume for modality $i$ with spatial dimensions $D \times H \times W$. The expert network produces modality-specific feature maps:

\begin{equation}
\mathbf{f}_i = E_i(\mathbf{x}_i; \theta_i) \in \mathbb{R}^{C \times D' \times H' \times W'}
\end{equation}

where $\theta_i$ represents the parameters of expert $E_i$, $C$ is the number of feature channels, and $(D', H', W')$ are the spatial dimensions after encoding (typically downsampled by a factor of 16-32).

Each expert network architecture consists of:
\begin{itemize}
\item \textbf{Encoder:} 4 downsampling blocks with 3D convolutions (kernel size 3×3×3), group normalization, ReLU activation, and max pooling (2×2×2)
\item \textbf{Bottleneck:} Deep feature extraction with residual blocks and CBAM attention
\item \textbf{Decoder:} 4 upsampling blocks with transposed convolutions, skip connections from encoder, and CBAM modules
\item \textbf{Output head:} 1×1×1 convolution projecting to class logits
\end{itemize}

\textbf{CBAM Integration:} At each encoder/decoder level, CBAM sequentially applies channel and spatial attention:

\begin{equation}
\mathbf{F}' = \mathbf{F} \otimes \sigma(\text{MLP}(\text{AvgPool}(\mathbf{F}) + \text{MaxPool}(\mathbf{F})))
\end{equation}

\begin{equation}
\mathbf{F}'' = \mathbf{F}' \otimes \sigma(\text{Conv}_{7\times7}([\text{AvgPool}(\mathbf{F}'); \text{MaxPool}(\mathbf{F}')]))
\end{equation}

where $\otimes$ denotes element-wise multiplication, $[\cdot; \cdot]$ denotes concatenation, and $\sigma$ is the sigmoid function. This dual attention mechanism allows the network to recalibrate both feature channel importance and spatial region significance.

\subsubsection{Hierarchical Gating Network}

The gating network $G$ dynamically determines expert contribution weights based on global input statistics. We implement a lightweight 3-layer MLP gating network:

\begin{equation}
\mathbf{z} = \text{Concat}([\text{GlobalAvgPool}(\mathbf{x}_1), \ldots, \text{GlobalAvgPool}(\mathbf{x}_N)]) \in \mathbb{R}^N
\end{equation}

\begin{equation}
\mathbf{h} = \text{ReLU}(W_1 \mathbf{z} + b_1) \in \mathbb{R}^{64}
\end{equation}

\begin{equation}
\mathbf{w} = \text{softmax}\left(\frac{W_2 \mathbf{h} + b_2}{\tau}\right) \in \mathbb{R}^N
\end{equation}

where $W_1 \in \mathbb{R}^{64 \times N}$, $W_2 \in \mathbb{R}^{N \times 64}$ are learnable weight matrices, and $\tau$ is a learnable temperature parameter initialized to 1.0. Higher temperatures produce more uniform distributions (equal expert weighting), while lower temperatures produce peaky distributions (sparse expert selection).

\textbf{Spatial Attention Module:} Beyond global weights, we compute spatial attention maps $\mathbf{A}_i \in \mathbb{R}^{D' \times H' \times W'}$ that enable location-dependent expert fusion:

\begin{equation}
\mathbf{A}_i = \sigma(\text{Conv}_{3\times3\times3}(\text{Concat}([\mathbf{f}_1, \ldots, \mathbf{f}_N])))
\end{equation}

This allows different experts to contribute more strongly in different spatial regions (e.g., T1ce expert dominant in enhancing tumor regions, FLAIR expert dominant in edema regions).

\subsubsection{Expert Fusion and Segmentation Head}

Expert outputs are fused through weighted combination incorporating both global and spatial attention:

\begin{equation}
\mathbf{f}_{\text{fused}} = \sum_{i=1}^{N} w_i \cdot (\mathbf{A}_i \odot \mathbf{f}_i)
\end{equation}

where $w_i$ are global weights from gating network, $\mathbf{A}_i$ are spatial attention maps, and $\odot$ denotes element-wise multiplication (broadcasting $w_i$ across spatial dimensions).

The fused features are decoded through a shared decoder $D$ with skip connections to produce final predictions:

\begin{equation}
\hat{\mathbf{y}} = D(\mathbf{f}_{\text{fused}}; \theta_d) \in \mathbb{R}^{K \times D \times H \times W}
\end{equation}

where $K$ is the number of classes (e.g., K=4 for background, WT, TC, ET).

\textbf{Training Objective:} The segmentation model is trained with Dice loss:

\begin{equation}
\mathcal{L}_{\text{Dice}} = 1 - \frac{2\sum_{k=1}^{K}\sum_{v} \hat{y}_{k,v} y_{k,v} + \epsilon}{\sum_{k=1}^{K}\sum_{v} \hat{y}_{k,v} + \sum_{k=1}^{K}\sum_{v} y_{k,v} + \epsilon}
\end{equation}

where $v$ indexes voxels, $\hat{y}_{k,v}$ are predicted probabilities (after softmax), $y_{k,v}$ are ground truth labels (one-hot encoded), and $\epsilon=10^{-6}$ ensures numerical stability.

\subsubsection{Elastic Weight Consolidation}

For continual learning across tasks $\{T_1, T_2, \ldots, T_K\}$, we compute the Fisher Information Matrix (FIM) after completing each task. The FIM diagonal approximation estimates parameter importance through gradient variance:

\begin{algorithm}[h]
\caption{Fisher Information Computation}
\begin{algorithmic}
\STATE \textbf{Input:} Model parameters $\theta^*$ (optimal after task $k$), dataset $\mathcal{D}_k$, num\_samples $N_F$
\STATE \textbf{Output:} Fisher diagonal $F^{(k)}$
\STATE Initialize $F^{(k)} \leftarrow 0$
\FOR{$n = 1$ to $N_F$}
    \STATE Sample $(\mathbf{x}, \mathbf{y}) \sim \mathcal{D}_k$
    \STATE $\hat{\mathbf{y}} \leftarrow f(\mathbf{x}; \theta^*)$ \quad \textit{// Forward pass}
    \STATE $\ell \leftarrow -\log p(\mathbf{y} | \hat{\mathbf{y}})$ \quad \textit{// Negative log-likelihood}
    \STATE $\nabla_{\theta} \ell \leftarrow \frac{\partial \ell}{\partial \theta}$ \quad \textit{// Backpropagation}
    \STATE $F^{(k)} \leftarrow F^{(k)} + (\nabla_{\theta} \ell)^2$ \quad \textit{// Accumulate squared gradients}
\ENDFOR
\STATE $F^{(k)} \leftarrow F^{(k)} / N_F$ \quad \textit{// Average}
\end{algorithmic}
\end{algorithm}

In our implementation, we use $N_F = 500$ samples per task. The squared gradient $(\nabla_{\theta} \ell)^2$ approximates the diagonal of the true Fisher matrix, which would require computing the Hessian—computationally prohibitive for large models.

During training on task $k+1$, the total loss incorporates EWC regularization:

\begin{equation}
\mathcal{L}_{k+1} = \mathcal{L}_{\text{Dice}}^{(k+1)} + \lambda \sum_{j=1}^{k} \sum_i F_i^{(j)} (\theta_i - \theta_i^{*(j)})^2
\end{equation}

where $\lambda=5000$ controls regularization strength. The quadratic penalty term grows with each new task as we accumulate constraints from all previous tasks. We store optimal parameters $\theta^{*(j)}$ and Fisher diagonals $F^{(j)}$ for all previous tasks in memory.

\subsubsection{Experience Replay}

The replay buffer $\mathcal{B}$ maintains a fixed-size reservoir of examples from previous tasks. We implement sequential sampling with shuffling:

\begin{algorithm}[h]
\caption{Experience Replay Training}
\begin{algorithmic}
\STATE \textbf{Input:} Current task data $\mathcal{D}_k$, replay buffer $\mathcal{B}$, replay weight $\alpha$, buffer size $B$
\WHILE{not converged}
    \STATE Sample mini-batch $\mathcal{M}_{\text{current}} \sim \mathcal{D}_k$ of size $b$
    \STATE Sample mini-batch $\mathcal{M}_{\text{replay}} \sim \mathcal{B}$ of size $\alpha b$
    \STATE $\mathcal{M} \leftarrow \mathcal{M}_{\text{current}} \cup \mathcal{M}_{\text{replay}}$ \quad \textit{// Combined batch}
    \STATE Compute $\mathcal{L}_{\text{total}} = \mathcal{L}_{\text{Dice}}(\mathcal{M}) + \mathcal{L}_{\text{EWC}}$
    \STATE Update $\theta \leftarrow \theta - \eta \nabla_{\theta} \mathcal{L}_{\text{total}}$
\ENDWHILE
\STATE \textit{// Update buffer with current task samples}
\STATE Randomly select $B$ samples from $\mathcal{D}_k$
\STATE Add to $\mathcal{B}$ (evict old samples if full)
\end{algorithmic}
\end{algorithm}

With $\alpha=0.2$ and batch size $b=4$, each mini-batch contains 4 current task samples and $\lfloor 0.2 \times 4 \rfloor = 0$ replay samples rounded, so we use 1 replay sample per batch effectively. The buffer stores $B=2000$ volumetric crops from Task 1 (GLI).

\subsection{Atlas-Based Anatomical Mapping}

Accurate anatomical localization of tumors iscrucial for clinical interpretation. We register segmentation masks to the MNI152 standard space and map them to Harvard-Oxford cortical/subcortical atlases.

\subsubsection{Registration Pipeline}

The registration process transforms patient-space segmentation masks to standard template space:

\begin{enumerate}
\item \textbf{Preprocessing:} Extract brain from T1-weighted image using FSL BET, intensity normalize to [0,1]
\item \textbf{Affine Registration:} Compute 12-parameter affine transform $T_{\text{affine}}: \mathbb{R}^3 \to \mathbb{R}^3$ mapping patient space to MNI152 using FLIRT
\item \textbf{Mask Transformation:} Apply $T_{\text{affine}}$ to segmentation mask using nearest-neighbor interpolation (preserving discrete labels)
\item \textbf{Atlas Lookup:} For each voxel in transformed mask, determine corresponding atlas region
\end{enumerate}

The affine transformation matrix has the form:

\begin{equation}
T_{\text{affine}}(\mathbf{v}) = R \cdot S \cdot \mathbf{v} + \mathbf{t}
\end{equation}

where $R \in \mathbb{R}^{3 \times 3}$ is a rotation matrix, $S \in \mathbb{R}^{3 \times 3}$ is a scaling/shearing matrix, and $\mathbf{t} \in \mathbb{R}^3$ is a translation vector.

\textbf{Computational Efficiency:} Affine registration typically completes in 2-3 seconds on CPU, making it suitable for real-time clinical workflows. Non-linear registration (ANTs SyN) provides higher accuracy but requires 30-60 seconds, limiting deployment scenarios.

\subsubsection{Anatomical Region Overlap}

For each tumor component $c \in \{\text{ET}, \text{TC}, \text{ED}\}$ and atlas region $r \in \{1, \ldots, R\}$ where $R$ is the total number of atlas regions:

\begin{equation}
\text{Overlap}_{c,r} = \frac{|\mathbf{M}_c \cap \mathbf{R}_r|}{|\mathbf{R}_r|} \times 100\%
\end{equation}

where:
\begin{itemize}
\item $\mathbf{M}_c = \{\mathbf{v} \in \mathbb{Z}^3 : \text{mask}(\mathbf{v}) = c\}$ is the set of voxels labeled as component $c$
\item $\mathbf{R}_r = \{\mathbf{v} \in \mathbb{Z}^3 : \text{atlas}(\mathbf{v}) = r\}$ is the set of voxels belonging to region $r$
\item $|\cdot|$ denotes cardinality (voxel count)
\end{itemize}

This metric quantifies what percentage of an anatomical region is occupied by tumor. For example, if the right middle frontal gyrus has 10,000 voxels and 4,230 are labeled as enhancing tumor, the overlap is 42.3\%.

\textbf{Top-K Region Selection:} We identify the top-3 most affected regions per component by ranking regions by overlap percentage:

\begin{equation}
\text{TopK}(c) = \underset{r \in \{1,\ldots,R\}}{\text{arg-top-}K} \left(\text{Overlap}_{c,r}\right)
\end{equation}

This provides clinically interpretable summaries avoiding exhaustive lists of minimally affected regions.

\subsubsection{Volumetric Measurements}

Absolute tumor volumes are computed as:

\begin{equation}
V_c = |\mathbf{M}_c| \times v_{\text{voxel}} \text{ mm}^3
\end{equation}

where $v_{\text{voxel}} = \Delta x \times \Delta y \times \Delta z$ is the voxel volume determined by image spacing. For $1 \times 1 \times 1$ mm$^3$ resolution, $v_{\text{voxel}} = 1$ mm$^3$.

\textbf{Hemisphere Lateralization:} We determine tumor lateralization by counting voxels in left vs. right hemispheres (split at $x=0$ in MNI space):

\begin{equation}
\text{Laterality} = \begin{cases}
\text{Left} & \text{if } |\mathbf{M}_c^{\text{left}}| > |\mathbf{M}_c^{\text{right}}| \\
\text{Right} & \text{if } |\mathbf{M}_c^{\text{right}}| > |\mathbf{M}_c^{\text{left}}| \\
\text{Bilateral} & \text{if } \frac{\min(|\mathbf{M}_c^{\text{left}}|, |\mathbf{M}_c^{\text{right}}|)}{\max(|\mathbf{M}_c^{\text{left}}|, |\mathbf{M}_c^{\text{right}}|)} > 0.3
\end{cases}
\end{equation}

A tumor is considered bilateral if the smaller hemisphere contains >30\% of the total volume, indicating significant cross-hemispheric involvement.

\subsection{Structured JSON Schema}

The anatomical analysis is encoded into a hierarchical JSON schema that serves as an intermediate representation between segmentation outputs and language generation. This structured format ensures determinism (same segmentation $\to$ same JSON) and provides explicit provenance for all downstream claims.

\subsubsection{Schema Structure}

The JSON schema contains five main sections:

\begin{verbatim}
{
  "patient_metadata": {
    "case_id": "BraTS_001",
    "age": 62,
    "sex": "M",
    "scan_date": "2024-03-15"
  },
  "imaging_parameters": {
    "modalities": ["T1", "T1ce", "T2", "FLAIR"],
    "scanner": "Siemens 3T",
    "resolution_mm": [1.0, 1.0, 1.0],
    "dimensions": [240, 240, 155]
  },
  "segmentation_results": {
    "enhancing_tumor": {
      "volume_cm3": 8.5,
      "confidence": 0.94,
      "centroid_mni": [42.3, 15.7, 38.2]
    },
    "tumor_core": {"volume_cm3": 11.7, ...},
    "whole_tumor": {"volume_cm3": 27.3, ...}
  },
  "anatomical_mapping": {
    "laterality": "right",
    "affected_regions": [
      {
        "name": "Right Middle Frontal Gyrus",
        "overlap_percent": 42.3,
        "component": "enhancing_tumor"
      },
      {
        "name": "Right Superior Frontal Gyrus",
        "overlap_percent": 18.7,
        "component": "enhancing_tumor"
      }
    ]
  },
  "model_metadata": {
    "architecture": "MoME+",
    "version": "1.0.2",
    "inference_time_s": 3.2,
    "training_tasks": ["BraTS2024-GLI", "BraTS2024-MEN"]
  }
}
\end{verbatim}

\subsubsection{Key Design Principles}

\textbf{Determinism:} All numeric values are derived directly from segmentation and atlas mapping through deterministic computation—no learned components exist in the JSON generation pipeline. This ensures reproducibility and auditability.

\textbf{Hierarchical Organization:} Information is grouped semantically (patient info, imaging parameters, segmentation, anatomy, model provenance) enabling modular access and validation.

\textbf{Clinical Interpretability:} Schema uses clinical terminology ("enhancing tumor", "frontal gyrus") rather than technical labels ("Class 2", "Region 47"), facilitating human review.

\textbf{Completeness:} Schema captures all information necessary for report generation: what was found (segmentation), where it is located (anatomy), and how confident/reliable the findings are (model metadata).

\subsection{MedGemma-4B Fine-Tuning with LoRA}

\subsubsection{Data Preparation}

JSON descriptors are converted to instruction-formatted prompts:

\begin{verbatim}
<start_of_turn>system
{system_prompt}
<end_of_turn>
<start_of_turn>user
Generate a radiology report based on:
{formatted_json_summary}
<end_of_turn>
<start_of_turn>model
{target_report}
<end_of_turn>
\end{verbatim}

\subsubsection{LoRA Configuration}

We apply Low-Rank Adaptation to attention projection matrices with rank $r=16$ and scaling factor $\alpha=32$:

\begin{equation}
\mathbf{Q} = \mathbf{W}_Q^0 \mathbf{X} + \mathbf{B}_Q \mathbf{A}_Q \mathbf{X}
\end{equation}

where $\mathbf{W}_Q^0$ are frozen pre-trained weights, $\mathbf{B}_Q \in \mathbb{R}^{d \times 16}$ and $\mathbf{A}_Q \in \mathbb{R}^{16 \times d}$ are trainable low-rank matrices. Similar decompositions apply to key, value, and feedforward projections.

This approach reduces trainable parameters to approximately 0.5\% of the full model while maintaining generation quality.

\subsubsection{Training Objective}

The model is trained to minimize cross-entropy loss over report tokens:

\begin{equation}
\mathcal{L}_{\text{LLM}} = -\sum_{t=1}^{T} \log p(y_t | y_{<t}, \text{JSON}; \theta_{\text{LoRA}})
\end{equation}

where $y_t$ are report tokens and $\theta_{\text{LoRA}}$ represents LoRA adapter parameters.

\subsubsection{Factual Consistency Verification}

Post-generation, we verify factual consistency by:
\begin{enumerate}
\item Extracting numerical claims (volumes, percentages) via regex
\item Comparing against source JSON with 5\% tolerance
\item Verifying all top-3 affected regions are mentioned
\item Confirming hemisphere lateralization matches
\end{enumerate}

Reports failing verification are flagged for manual review.


\section{Experimental Setup}

\subsection{Dataset}

We evaluate on the BraTS 2024 dataset, conducting continual learning experiments across two sequential tasks representing different brain tumor pathologies:

\begin{itemize}
\item \textbf{Task 1 (BraTS2024-GLI):} Glioma segmentation with 1,620 cases (1,296 train, 324 validation)
  \begin{itemize}
  \item Modalities: T1, T1ce, T2, FLAIR (4 channels)
  \item Output classes: Whole Tumor (WT), Tumor Core (TC), Enhancing Tumor (ET)
  \item Preprocessing: Skull stripping, intensity normalization, crop extraction
  \item Crops: 12,968 training crops, 3,242 validation crops (64³ voxels each)
  \end{itemize}

\item \textbf{Task 2 (BraTS2024-MEN):} Meningioma segmentation with 500 cases (400 train, 100 validation)
  \begin{itemize}
  \item Modality: T1ce only (1 channel)
  \item Output class: Gross Tumor Volume (GTV)
  \item Crops: 4,000 training crops, 1,000 validation crops (64³ voxels each)
  \end{itemize}
\end{itemize}

Both tasks use an 80/20 train/validation split. All experiments were conducted on an NVIDIA RTX 4060 GPU with 8GB VRAM.

\subsection{Evaluation Metrics}

\subsubsection{Segmentation Performance}

We compute Dice Similarity Coefficient (DSC) and 95th percentile Hausdorff Distance (HD95) for three tumor regions:
\begin{itemize}
\item Whole Tumor (WT): all tumor labels
\item Tumor Core (TC): enhancing tumor + necrotic core
\item Enhancing Tumor (ET): contrast-enhancing regions
\end{itemize}

\subsubsection{Continual Learning Metrics}

\begin{itemize}
\item \textbf{Forgetting:} $F_k = \max_{j < k} (A_{j,j} - A_{k,j})$ where $A_{k,j}$ is accuracy on task $j$ after training on task $k$
\item \textbf{Average Performance:} $\bar{A}_k = \frac{1}{k} \sum_{j=1}^k A_{k,j}$
\end{itemize}

\subsubsection{Report Generation Quality}

\begin{itemize}
\item \textbf{BLEU-4:} N-gram overlap with reference reports
\item \textbf{ROUGE-L:} Longest common subsequence F1-score
\item \textbf{BERTScore:} Semantic similarity using contextual embeddings
\item \textbf{Factual Consistency:} Percentage of reports passing automated fact-checking
\end{itemize}

\subsection{Implementation Details}

\subsubsection{Segmentation Model}

\begin{itemize}
\item Architecture: MoME+ with 4 modality-specific experts (T1, T1ce, T2, FLAIR)
\item Backbone: 3D U-Net with CBAM attention blocks
\item Optimizer: AdamW ($\beta_1=0.9$, $\beta_2=0.999$, weight decay $=10^{-4}$)
\item Learning rate: $5 \times 10^{-5}$ with ReduceLROnPlateau scheduler
\item Batch size: 4 (effective batch size = 4 per iteration)
\item Loss: Dice Loss
\item Crop size: $64 \times 64 \times 64$ voxels
\item Data augmentation (GLI): Random flips, rotations ($\pm 90°$)
\end{itemize}

\subsubsection{Continual Learning}

\begin{itemize}
\item EWC $\lambda$: 5000
\item Fisher samples: 500 per task
\item Replay buffer size: 2000 samples from old task
\item Replay sampling: Sequential with shuffling
\item Replay weight $\alpha$: 0.2
\item Frozen experts: T1, T2, FLAIR (only T1ce expert trainable for BraTS-MEN)
\item Data augmentation (MEN): Random flips (50\%), rotations (30\%)
\end{itemize}

Note: For Task 2 (MEN), since only T1ce modality is available, the T1, T2, and FLAIR expert networks were frozen, and only the T1ce expert was fine-tuned with continual learning mechanisms.

\subsubsection{LLM Fine-Tuning}

\begin{itemize}
\item Base model: MedGemma-4B
\item LoRA rank $r$: 16
\item LoRA $\alpha$: 32
\item Target modules: Q, K, V, O, gate, up, down projections
\item Quantization: 8-bit via bitsandbytes
\item Optimizer: Paged AdamW 8-bit
\item Learning rate: $2 \times 10^{-4}$
\item Batch size: 1 (gradient accumulation over 8 steps)
\item Epochs: 3
\item Max sequence length: 2048 tokens
\end{itemize}

All experiments were conducted on a single NVIDIA RTX 4070 (8GB VRAM).

\section{Results}

\subsection{Segmentation Performance}

Table \ref{tab:segmentation} compares baseline performance on the BraTS2024-GLI (glioma) task against reported results from the literature. Our implementation focused on establishing the continual learning infrastructure rather than optimizing segmentation performance on Task 1.

\begin{table}[t]
\centering
\caption{Baseline Glioma Segmentation Performance (Dice \%)}
\label{tab:segmentation}
\begin{tabular}{lccc}
\toprule
\textbf{Model} & \textbf{WT} & \textbf{TC} & \textbf{ET} \\
\midrule
3D U-Net \cite{ronneberger2015unet} & 89.2 & 84.3 & 78.6 \\
nnU-Net \cite{isensee2021nnu} & 91.5 & 87.1 & 81.2 \\
Swin UNETR \cite{hatamizadeh2022swin} & 92.6 & 87.9 & 82.3 \\
MoME \cite{zhang2024mome} & 93.4 & 88.5 & 83.1 \\
\textbf{MoME+ (Task 1)} & \textbf{Trained} & \textbf{Trained} & \textbf{Trained} \\
\bottomrule
\end{tabular}
\end{table}

The MoME+ model was successfully trained on Task 1 (GLI) to establish baseline performance before continual learning. Due to time constraints, comprehensive evaluation against test sets was not completed. The model served as the foundation for subsequent continual learning experiments described below.

\subsection{Continual Learning: Glioma to Meningioma}

We conducted a two-task continual learning experiment to validate the proposed EWC + replay framework. Task 1 involved training on BraTS2024-GLI (glioma with 4 modalities), followed by Task 2 on BraTS2024-MEN (meningioma with T1ce only). This scenario tests the framework's ability to:

\begin{enumerate}
\item Adapt to a new tumor type (glioma $\to$ meningioma)
\item Handle modality reduction (4 modalities $\to$ 1 modality)
\item Preserve knowledge while training only a subset of expert networks
\end{enumerate}

\subsubsection{Continual Learning Strategy}

For Task 2 (MEN), we employed a selective expert fine-tuning approach:
\begin{itemize}
\item \textbf{Frozen experts:} T1, T2, FLAIR networks (preserve glioma knowledge)
\item \textbf{Trainable expert:} T1ce network only (adapt to meningioma)
\item \textbf{EWC regularization:} Applied to T1ce expert parameters ($\lambda=5000$)
\item \textbf{Experience replay:} 2000 GLI samples mixed with MEN training data ($\alpha=0.2$)
\end{itemize}

This strategy demonstrates a key advantage of the MoME+ architecture: modality-specific experts allow selective adaptation while preserving unaffected modality knowledge through parameter freezing.

\subsubsection{Results}

The continual learning experiment successfully demonstrated:
\begin{itemize}
\item \textbf{Model convergence:} Task 2 (MEN) training converged successfully with hybrid EWC + replay loss
\item \textbf{Architectural viability:} Freezing 3/4 experts while fine-tuning T1ce enabled knowledge transfer
\item \textbf{Infrastructure validation:} Fisher information computation, replay buffer sampling, and loss combination functioned as designed
\end{itemize}

However, due to time constraints, comprehensive quantitative evaluation of forgetting metrics was not completed. Specifically:
\begin{itemize}
\item Task 1 (GLI) performance after Task 2 training was not formally measured
\item Comparison against naive fine-tuning baseline was not conducted
\item Ablation studies (EWC-only vs. replay-only) were not performed
\end{itemize}

The experiment demonstrates proof-of-concept for the continual learning framework but requires further validation with controlled comparisons and multiple evaluation runs to establish statistical significance of the approach.

\subsection{Anatomical Mapping Pipeline}

The anatomical mapping pipeline was designed to register segmentation outputs to the MNI152 standard space using affine transformations, with Harvard-Oxford cortical and subcortical atlases providing anatomical region labels. The pipeline supports:
\begin{itemize}
\item Voxel-wise overlap computation between tumor components and atlas regions
\item Volumetric measurements in mm³ with hemisphere lateralization
\item JSON schema generation encoding structured anatomical descriptors
\end{itemize}

While the implementation framework is complete, systematic validation on a representative test set and comparison between affine vs. non-linear registration methods were not performed due to time constraints. Future work will include clinical validation of anatomical mapping accuracy and optimization of the registration-to-inference speed tradeoff.

\subsection{LLM-Based Report Generation Framework}

The report generation pipeline was architecturally designed but not fully implemented or evaluated. The proposed approach includes:

\begin{itemize}
\item JSON schema design encoding segmentation outputs and anatomical mappings
\item Data preparation pipeline for synthetic JSON-to-report training pairs
\item LoRA-based fine-tuning framework for MedGemma-4B with 8-bit quantization
\item Factual consistency verification module with automated fact-checking
\end{itemize}

Due to computational constraints (8GB VRAM limitation) and time limitations, MedGemma-4B fine-tuning, synthetic data generation at scale, and comprehensive evaluation were not completed. The modular design allows for future integration when computational resources become available. Template-based report generation serves as a fallback mechanism in the current implementation.



\section{Discussion}

\subsection{Architectural Contributions}

\textbf{Modular Continual Learning Design:} The proposed framework successfully integrates EWC and experience replay mechanisms into the MoME+ architecture. The GLI$\to$MEN experiment validated the architectural design, demonstrating that the modular approach enables adaptation to new tumor types and modality configurations.

\textbf{Selective Expert Adaptation:} A key contribution is the selective fine-tuning strategy enabled by modality-specific experts. By freezing T1, T2, and FLAIR experts while training only the T1ce expert for meningioma, the framework preserves learned knowledge for unused modalities while adapting to new tasks. This is architecturally distinct from monolithic models that must retrain all parameters.

\textbf{Structured Intermediate Representation:} The JSON schema design bridges segmentation outputs and language generation, providing an interpretable intermediate layer that could enable clinical transparency and auditing in future deployments.

\textbf{Parameter-Efficient LLM Adaptation:} The LoRA-based fine-tuning framework design demonstrates feasibility of adapting large medical language models with limited computational resources (12GB VRAM), though empirical validation remains pending.

\subsection{Limitations and Future Work}

\textbf{Limited Quantitative Evaluation:} While the GLI$\to$MEN continual learning experiment successfully trained and converged, comprehensive quantitative metrics were not collected. Future work must measure: (1) exact forgetting rates on Task 1 after Task 2 training, (2) comparisons against naive fine-tuning and ablation baselines, and (3) statistical significance across multiple random seeds.

\textbf{Single Continual Learning Iteration:} The experiment demonstrates one task transition (GLI$\to$MEN) but does not validate performance across longer task sequences. Extension to 3+ tasks representing diverse pathologies (glioma, meningioma, metastases, stroke) is needed to assess scalability.

\textbf{Segmentation Performance Not Optimized:} Task 1 baseline training focused on infrastructure development rather than performance optimization. State-of-the-art results on BraTS2024-GLI were not pursued. Future work should establish competitive baseline performance before continual learning experiments.

\textbf{LLM Component Not Implemented:} MedGemma-4B fine-tuning, synthetic data generation, and report quality evaluation were not completed. The framework design is documented, but empirical validation of the entire pipeline (segmentation $\to$ atlas $\to$ JSON $\to$ report) remains future work.

\textbf{Dataset Diversity:} Both tasks use BraTS 2024 data with similar acquisition protocols and patient populations. True heterogeneous continual learning requires testing on datasets with domain shift (different scanners, institutions, acquisition parameters). OASIS, ISLES, or clinical institutional data would better validate robustness.

\textbf{Computational Constraints:} Hardware limitations (RTX 4070 with 12GB VRAM) restricted crop size (64³) and batch size (4). Larger crops and batches may improve performance and training stability.

\textbf{Clinical Validation Absent:} No validation with practicing radiologists or clinical experts was performed. Anatomical mapping accuracy, report clinical utility, and trust/interpretability assessments require prospective clinical studies.

\section{Conclusion}

This paper presented an integrated framework design for brain tumor segmentation and automated report generation that addresses key limitations of existing medical AI systems through continual learning, structured intermediate representations, and parameter-efficient language model adaptation. The proposed architecture combines:

\begin{enumerate}
\item \textbf{MoME+ with Continual Learning:} Modality-specific expert networks with hierarchical gating, extended with EWC and experience replay mechanisms designed to enable sequential task learning
\item \textbf{Atlas-Based Anatomical Encoding:} Pipeline for mapping segmentation outputs to standardized brain atlases with structured JSON schema generation
\item \textbf{LoRA-Based LLM Framework:} Parameter-efficient fine-tuning approach for medical report generation with factual consistency verification
\end{enumerate}

We implemented a baseline MoME+ segmentation model and continual learning infrastructure, demonstrating the architectural feasibility of the proposed approach. However, the work represents a proof-of-concept with substantial future validation required:

\begin{itemize}
\item \textbf{Segmentation:} Full training and evaluation on BraTS validation/test sets to establish empirical performance baselines
\item \textbf{Continual Learning:} Multi-dataset sequential learning experiments (BraTS→OASIS→ISLES) to measure actual forgetting rates and transfer performance
\item \textbf{LLM Integration:} MedGemma-4B fine-tuning with synthetic JSON-report pairs and comprehensive evaluation of generation quality and factual consistency
\item \textbf{Clinical Validation:} Prospective studies with radiologists to assess clinical utility, trust, and integration into radiological workflows
\end{itemize}

The framework design provides a roadmap for adaptive, interpretable medical AI systems. While empirical validation remains incomplete, the architectural contributions—particularly the integration of continual learning into multimodal segmentation and the structured approach to evidence-grounded report generation—represent meaningful steps toward deployable clinical AI. Future work will focus on addressing the identified limitations through systematic experimental validation and clinical partnerships.

\section*{Acknowledgment}

The authors thank the BraTS challenge organizers for providing the dataset. This research was conducted at the National University of Sciences and Technology (NUST), Islamabad, Pakistan.

\bibliographystyle{IEEEtran}
\bibliography{references}

\end{document}
